\documentclass{letter}
\usepackage{hyperref}
\signature{[Author Name]}
\address{[Affiliation]}
\begin{document}

\begin{letter}{Editor-in-Chief\\Advanced Engineering Informatics}

\opening{Dear Editor,}

Please find enclosed our manuscript ``Learning one magnetic model for all
winding connections: connection-invariant co-energy surrogates and
drive-cycle co-design of reconfigurable electric machines'', which we submit
for consideration as a research paper in Advanced Engineering Informatics.

Reconfigurable windings are an established route to wide-speed-range,
high-efficiency electrified powertrains, yet every study we are aware of
evaluates each candidate connection with its own dedicated FE model or
flux-map table, and schedules mode changes with fixed thresholds or greedy
rules. Our manuscript addresses both gaps with one idea: a winding
connection is a set of linear constraints on coil currents, so a single
connection-invariant surrogate---trained once on coil-level FE data---can
evaluate \emph{any} connection zero-shot. We think this fits the journal's
scope in three specific ways: (i) it is an engineering-informatics method
(learned surrogates embedded in design and operation workflows) rather than
a new machine or a new controller; (ii) it comes with a reproducible
benchmark---7 hold-out connections, three difficulty tiers, 9{,}276+
nonlinear FE samples, fixed splits and protocols---released with code and
data; and (iii) it delivers actionable engineering findings, including two
data-collection design principles (what \emph{kind} of current patterns to
sample, which matters more than how many), a quantified physical-consistency
advantage of co-energy-structured surrogates, an instructive negative result
on monotonicity-constrained architectures, and a drive-cycle dynamic-program
    that prices reconfiguration honestly (changeover losses, dwell constraints).
    An independent FEMM cross-check is included; it validates the coarse
    inductance and nonlinear flux trends while exposing a zero-shot torque
    failure for an untrained tap connection, so the manuscript's claims are
    explicitly scoped to what the evidence supports.

All results in the manuscript are reproducible from the released code and
data with one command per figure; the FE reference data, trained models, and
benchmark splits are included. The work has not been published elsewhere and
is not under consideration elsewhere. The authors declare no competing
interests.

We suggest the following potential reviewers with relevant expertise:
[3--6 suggested reviewers with affiliations and emails].

Thank you for your consideration.

\closing{Sincerely,}

\end{letter}
\end{document}
